\chapter*{Abkürzungsverzeichnis}
\addcontentsline{toc}{chapter}{Abkürzungsverzeichnis}
% \chapter* erzeugt eine Überschrift OHNE Nummerierung
% \addcontentsline sorgt dennoch dafür, dass sie im Inhaltsverzeichnis erscheint

% ============================================================
%  VARIANTE 1: Manuell (einfach, aber bei Änderungen aufwändig)
% ============================================================
% Einfach weitere Zeilen nach dem Muster "KÜRZEL & Bedeutung \\" ergänzen.

%\begin{longtable}{@{}p{3cm}p{12cm}@{}}
%API     & Application Programming Interface \\
%HSWT    & Hochschule Weihenstephan-Triesdorf \\
%KI      & Künstliche Intelligenz \\
%PDF     & Portable Document Format \\
% ... weitere Abkürzungen hier ergänzen
%\end{longtable}



% ============================================================
%  VARIANTE 2: Automatisiert mit dem acronym-Paket (empfohlen)
% ============================================================
% Vorteile:
%  - Jede Abkürzung nur einmal definieren
%  - Beim ersten Auftreten im Text wird automatisch "Langform (Kurzform)" ausgegeben
%  - Danach erscheint nur noch die Kurzform
%
% Schritt 1: Paket in main.tex laden:
%   \usepackage{acronym}
%
% Schritt 2: Abkürzungen hier definieren:
%
\begin{acronym}
    \acro{KI}{Künstliche Intelligenz}
    \acro{HSWT}{Hochschule Weihenstephan-Triesdorf}
    \acro{API}{Application Programming Interface}
\end{acronym}
%
% Schritt 3: Im Fließtext verwenden:
%   \ac{KI}   -> erstes Vorkommen: "Künstliche Intelligenz (KI)"
%             -> danach automatisch: "KI"
%
% Möchte man die Langform nochmal erzwingen: \acf{KI}
% Nur die Kurzform:                          \acs{KI}
% Nur die Langform:                          \acl{KI}
